\chapter{Discussion}
\label{ch:discussion}

\epigraph{The first principle is that you must not fool yourself---and you are the
easiest person to fool.}{Richard P. Feynman}

This chapter evaluates the thesis against its own questions and threats. It
revisits the research questions with bounded verdicts (\cref{sec:disc-rqs}), states
the threats to validity in full (\cref{sec:disc-threats}), enumerates limitations
and open items (\cref{sec:disc-limits}), and considers ethical and societal
consequences (\cref{sec:disc-ethics}).

\section{The Research Questions, Revisited}
\label{sec:disc-rqs}

\paragraph{RQ1 (shared structure; one substrate) --- \admitted\ analytically,
\candidate\ empirically.} That LSP, MCP, and A2A instantiate one decoupling is
\admitted: they share JSON-RPC with capability negotiation
\citep{lsp-spec-318,mcp-spec-2025,a2a-spec}, acknowledged lineage
\citep{mcp-spec-2025,a2a-mcp}, and converging governance
\citep{lf-a2a-press,lf-aaif-press}, and the Decoupling Theorem (\cref{thm:decoupling})
formalises the common mechanism. That their \emph{composition} forms a single
operational substrate is \candidate: it is demonstrated by construction
(\cref{ch:fusion}) and instantiated partly in the artifact, but not yet at
ecosystem scale.

\paragraph{RQ2 (when a domain is a language) --- \admitted\ as a criterion,
\candidate\ in demonstration.} The criterion---a domain is servable when it admits
an object-centric event-log abstraction (\cref{def:eventlog})---is stated and
defended \admitted. Its application across six domains (\cref{ch:all-language}) is
\candidate: the demonstrations are illustrative scenarios, not deployed systems
(\cref{sec:m-eval}).

\paragraph{RQ3 (process mining for all language) --- \candidate.} The thesis shows,
by construction and through the artifact's own OCEL binding (\cref{sec:art-ocel}),
that conformance checking applies beyond business process, and that the
object-centric log is the bridge (\cref{sec:al-bridge}). Field validation across
domains is \openstat, and the strongest evidence---the artifact reproducing
standard object-centric conformance on its calibration domain---is partial.

\paragraph{RQ4 (a trustworthy law-state runtime) --- \admitted\ for the four
properties, with \partialstat/\openstat\ sub-items.} The artifact realises the four
properties of \cref{def:lawstate}: object-centric observation, a conformance
surface, protocol exposure, and receipt-bound admission, each traced to source
(\cref{ch:artifact}). Three-valued non-collapse is \admitted\ (\cref{thm:noncollapse},
\cref{lst:cv}); read-only observation is \admitted\ by construction; receipt
chaining is \admitted. Sub-items remain \openstat, notably subagent gate
propagation, which the artifact itself records as \openstat\ rather than masking.

\paragraph{RQ5 (phase transition; 2030) --- \admitted\ as a model, \forecast\ as a
prediction.} The phase-transition account is \admitted\ as a model with a stated
mechanism (\cref{thm:decoupling} gearing $\times$ \cref{prop:landau-jump} threshold)
and a rigorous metaphor (\cref{thm:clausius}). Its application to 2030 is a
\forecast\ (\cref{ch:vision}), bounded by \cref{prin:plateau}.

\section{Threats to Validity}
\label{sec:disc-threats}

\paragraph{Construct validity.} The chief threat is the metaphor. If ``phase
transition'' is read literally---as a claim that information obeys
thermodynamics---the construct is invalid. \Cref{sec:fw-phase} confines the claim to
\emph{order type} (a discontinuity) established by an independent mechanism, and
\cref{thm:clausius} grounds only the metaphor's internal consistency, not a physical
identity. A second construct threat is the word ``conformance'' itself: the thesis
uses it in van der Aalst's technical sense, and its extension to, say, legal or
musical ``law'' is an analogy whose precision varies by domain
(\cref{ch:all-language}).

\paragraph{Internal validity.} The artifact is analysed from source and
specification, not from a live deployment exercised across all six domains
(\cref{sec:m-artifact}). Claims about its runtime behaviour are therefore warranted
to the extent that source implies behaviour---strong for type-level invariants
(\cref{thm:noncollapse} mirrors a compiled bitmask invariant), weaker for emergent
or concurrent behaviour. The sibling \code{wasm4pm} engine and its type authority
were not present in the study environment and are characterised from the artifact's
interfaces; their internals are \unknownstat.

\paragraph{External validity.} Generalising from one artifact and six illustrative
domains to ``all language'' is the largest threat. The event-log criterion
(\cref{def:eventlog}) is the explicit boundary: domains whose use cannot be observed
as timed events over typed objects fall outside the claim, and the thesis does not
assert universality beyond that criterion. The six domains were chosen to span text,
formal systems, sequence, and time-structured notation, but they are not a random
sample.

\paragraph{Reliability and reproducibility.} Primary claims are traced to primary
sources and to file-and-line locations so a later reader can re-examine them. The
recent-advances citations (\cref{sec:recent}) were machine-harvested with a
fifteen-item hand-verified spot-check; as living preprints they are \candidate\ and
should be re-verified. The forecasts (\cref{sec:vis-forecasts}) are third-party and
time-stamped.

\section{Limitations and Open Items}
\label{sec:disc-limits}

The thesis is conceptual and design-scientific, and three limitations follow.
First, no longitudinal deployment substantiates the 2030 expansion; the prediction
is a \forecast\ and may be falsified by a plateau that does not end
(\cref{prin:plateau}). Second, the artifact carries its own \openstat\ items---most
clearly that pre-tool-use gate enforcement does not cross subagent session
boundaries---and the thesis inherits them rather than resolving them. Third, the
mathematical pillars (\cref{ch:math}) derive the thesis's constructs but make no new
mathematical claim; their contribution is rigour and unification, not novelty in
algebra, analysis, or geometry. Each limitation is recorded here so that no later
sentence is read as more than it is.

\section{Ethical and Societal Considerations}
\label{sec:disc-ethics}

A substrate that lets a single conformance verdict travel across every agent
boundary is a concentration of authority, and concentrations invite three concerns.
\emph{Capture}: whoever defines the law axes defines admissibility; neutral
governance (\cref{sec:vis-governance}) mitigates but does not eliminate this, and
the thesis takes no position that a foundation cannot be captured. \emph{Over-trust}:
a receipt attests that a check ran, not that the check was right; ``proof-of-guardrail''
must be trusted only as far as its analysis warrants \citep{jin26guardrail}, and the
oversight stratum inherits the finite human-attention capacity that can make more
escalation less safe \citep{turan26}. \emph{Conformance theatre}: the easiest way to
pass a conformance gate is to weaken the law, and a system that rewards green
verdicts will be pressured to do so---which is precisely why the artifact forbids
victory language and requires receipts with negative controls
(\cref{sec:art-canary}). The dual-use character is real: the same machinery that
holds an agent accountable can surveil a human author whose prose, law, or code is
being ``conformance-checked''. The thesis's read-only, three-valued, evidence-bound
posture is a partial safeguard---it refuses silent mutation and refuses to call the
unchecked compliant---but it is not a complete answer, and presenting it as one would
violate the discipline the thesis adopts.
